\begin{table}[!htbp]\small\setlength{\tabcolsep}{4pt}
\caption{Composantes du nouvel impôt sur le revenu (Md€ de 2025, 2032 et à terme)}\label{tab:t2}
\begin{tabularx}{\linewidth}{@{}Xrc@{}}
\toprule
\textbf{Composante} & \textbf{Md€} & \textbf{Statut} \\
\midrule
Suppression des niches (réductions et crédits, exonérations, demi-parts) & +19,2 & DED \\
Pensions imposées comme les salaires (abattement, CSG, âge) & +15,6 & DED \\
Forfait de 500 € au lieu de l'abattement de 10 % et des frais réels & +25,7 & DED \\
Fin de la CSG déductible & +24,1 & DED \\
Impôt individuel, fin du quotient conjugal & +12,1 & DED \\
Fin du quotient familial & +16,2 & DED \\
Franchise de 4\,700 € par enfant & $-$15,1 & HYP \\
Baisse de tous les taux (nouveau barème) & $-$97,1 & CAL \\
Minimum vieillesse, contribution sur les hauts revenus & $-$0,8 & DED \\
\midrule \textbf{Solde public} & \textbf{0,0} & \\
\bottomrule
\end{tabularx}
\par\vspace{2pt}{\raggedright\scriptsize\color{gris}Source : modèle de foyers calé sur les statistiques de l'impôt sur le revenu 2024 (DGFiP) ; OpenFisca-France ; coûts des niches : Évaluation des voies et moyens (PLF 2026). Décomposition séquentielle avant réaction des comportements : l'ordre change la taille de chaque ligne, pas le total. Statut : DED, déduit d'une donnée observée ; CAL, calé (ici sur la neutralité par rapport au droit de 2032) ; HYP, paramètre choisi.\par}
\end{table}
